\section{Reinforcement Learning Foundations}
\label{sec:reinforcement_learning_foundations}

Reinforcement learning provides a mathematical framework for learning behavior through interaction with an environment.
The standard formulation is an episodic Markov decision process (MDP), defined by an observation or state space $\mathcal{S}$, an action space $\mathcal{A}$, a transition distribution $P(s_{t+1}\mid s_t,a_t)$, a reward function $r(s_t,a_t,s_{t+1})$, a discount factor $\gamma\in[0,1)$, and an episode horizon.
At each time step $t$, the agent observes $s_t$, samples or selects an action $a_t$, receives reward $r_t$, and transitions to $s_{t+1}$.
The goal is to learn a policy $\pi_\theta(a\mid s)$, parameterized by $\theta$, that maximizes expected discounted return,
\begin{equation}
    J(\theta)=\mathbb{E}_{\pi_\theta}\left[\sum_{t=0}^{T-1}\gamma^t r_t\right].
\end{equation}
This interaction-based formulation makes reinforcement learning distinct from supervised learning because the data distribution depends on the policy being trained.
Consequently, the agent must balance exploitation of currently rewarding actions with exploration of actions whose long-term consequences are uncertain.

Value functions summarize the long-term utility of states and actions.
The state-value function of a policy is
\begin{equation}
    V^\pi(s)=\mathbb{E}_{\pi}\left[\sum_{k=0}^{T-t-1}\gamma^k r_{t+k}\mid s_t=s\right],
\end{equation}
while the action-value function is
\begin{equation}
    Q^\pi(s,a)=\mathbb{E}_{\pi}\left[\sum_{k=0}^{T-t-1}\gamma^k r_{t+k}\mid s_t=s, a_t=a\right].
\end{equation}
The difference $A^\pi(s,a)=Q^\pi(s,a)-V^\pi(s)$ is the advantage function, which measures whether an action is better or worse than the policy's expected behavior at the same state.
Advantage estimates are central in modern policy-gradient algorithms because they reduce variance while preserving the direction of policy improvement.

Deep reinforcement learning extends this framework by using neural networks to approximate policies, value functions, and auxiliary models.
This approximation is necessary in continuous-control and high-dimensional settings where tabular representations are infeasible.
However, function approximation also makes exploration more difficult.
A neural policy may repeatedly visit a narrow subset of the state space if initial rewards favor locally stable behavior, and a value function may generalize poorly when the agent has not collected diverse experience.
Sparse-reward environments make the problem more severe because most early transitions may have identical or nearly identical external rewards.

Reward design is therefore a central issue in reinforcement learning.
Potential-based reward shaping can preserve optimal policies when it follows a restricted form, but arbitrary shaping terms may change the task objective \,\cite{ng-1999}.
Intrinsic-motivation methods intentionally add such auxiliary shaping signals to improve exploration, especially when the original reward does not provide enough feedback during early training.
These methods should therefore be evaluated not only by whether they increase the reward used for optimization, but also by whether they improve the extrinsic return of the learned policy.

The exploration problem addressed in this study lies at the intersection of sparse feedback, function approximation, and reward augmentation.
An effective intrinsic reward should encourage the agent to collect informative transitions without permanently distracting policy optimization from the external task.
This requirement motivates related work on novelty, prediction error, impact, and learning progress, each of which defines a different proxy for the value of experience.